% GENERATED FILE. Edit canonical lessons, then run npm run generate:books.
% canonical-source-hash: fnv1a64:e500257d5d23b883
% canonical-lessons: FR-C417-douanier, FR-C417-declarer, FR-C417-portique, FR-C417-fouiller, FR-C417-transiter

\chapter{À la douane}
\label{ch:fr-a-la-douane}
\hlchaptermodality{417}

\begin{chapteropening}
\textbf{By the end of this chapter:} \emph{I can go through security and customs: the customs officer, declaring goods, the scanner gate, searching bags and changing planes.}

Complete the last lesson of chapter 417: I can go through security and customs: the customs officer, declaring goods, the scanner gate, searching bags and changing planes.
\end{chapteropening}

\section[le douanier]{le douanier — the customs officer}

\label{lesson:FR-C417-douanier}



\begin{quote}
Before the new one: say the French for the hold, then the French for the boarding.
\end{quote}

\subsection*{You'll want to know: le douanier}
\textbf{le douanier} — \textquotedblleft{}the customs officer\textquotedblright{}.

\begin{quote}
Le douanier a demandé ce qu'il y avait dans la boîte.
\end{quote}

\begin{quote}
\emph{The customs officer asked what was in the box.}
\end{quote}

Built on \emph{la douane} 'customs'.

\begin{cousinweb}
\emph{douanier} is built on \emph{douane} 'customs', which is usually traced through Old Italian \emph{doana} to Arabic \emph{dīwān} 'register, office'. EtymDB does not list the word.
\end{cousinweb}

\begin{grammarlens}[title={What you've built}]
One more word for telling, explaining and describing at B1.
\end{grammarlens}

\subsection*{Your turn}
\begin{itemize}
\raggedright
  \item \emph{Say it:} \emph{le douanier}
  \item \emph{Say it:} \emph{le douanier}, once more
  \item \emph{Say it:} a sentence of your own with \emph{le douanier}
  \item \emph{Recall:} say the French for the hold, then the French for the boarding, then say \emph{le douanier} again
\end{itemize}

\begin{tcolorbox}[breakable,colback=teal!4,colframe=teal!35!black,title={Before you move on}]
What does le douanier mean? (\textquotedblleft{}The customs officer\textquotedblright{}.) Say it once more.
\end{tcolorbox}
\section[déclarer]{déclarer — to declare}

\label{lesson:FR-C417-declarer}



\begin{quote}
Before the new one: say the French for the boarding, then the French for the customs officer.
\end{quote}

\subsection*{You'll want to know: déclarer}
\textbf{déclarer} — \textquotedblleft{}to declare\textquotedblright{}.

\begin{quote}
Elle n'a rien à déclarer.
\end{quote}

\begin{quote}
\emph{She has nothing to declare.}
\end{quote}

Also \emph{déclarer ses revenus} — to file a tax return.

\begin{cousinweb}
\emph{déclarer} is a learned borrowing from Latin \emph{dēclārāre} 'to make clear', from \emph{clārus} 'clear' (EtymDB).
\end{cousinweb}

\begin{grammarlens}[title={What you've built}]
One more word for telling, explaining and describing at B1.
\end{grammarlens}

\subsection*{Your turn}
\begin{itemize}
\raggedright
  \item \emph{Say it:} \emph{déclarer}
  \item \emph{Say it:} \emph{déclarer}, once more
  \item \emph{Say it:} a sentence of your own with \emph{déclarer}
  \item \emph{Recall:} say the French for the boarding, then the French for the customs officer, then say \emph{déclarer} again
\end{itemize}

\begin{tcolorbox}[breakable,colback=teal!4,colframe=teal!35!black,title={Before you move on}]
What does déclarer mean? (\textquotedblleft{}To declare\textquotedblright{}.) Say it once more.
\end{tcolorbox}
\section[le portique]{le portique — the scanner gate (at security); the portico}

\label{lesson:FR-C417-portique}



\begin{quote}
Before the new one: say the French for the customs officer, then the French for to declare.
\end{quote}

\subsection*{You'll want to know: le portique}
\textbf{le portique} — \textquotedblleft{}the scanner gate (at security); the portico\textquotedblright{}.

\begin{quote}
Au contrôle, le portique a sonné à cause de sa ceinture.
\end{quote}

\begin{quote}
\emph{At security the gate beeped because of his belt.}
\end{quote}

\emph{Un portique de sécurité} — a walk-through scanner.

\begin{cousinweb}
\emph{portique} is a learned borrowing from Latin \emph{porticus} 'colonnade', built on \emph{porta} 'gate' (EtymDB).
\end{cousinweb}

\begin{grammarlens}[title={What you've built}]
One more word for telling, explaining and describing at B1.
\end{grammarlens}

\subsection*{Your turn}
\begin{itemize}
\raggedright
  \item \emph{Say it:} \emph{le portique}
  \item \emph{Say it:} \emph{le portique}, once more
  \item \emph{Say it:} a sentence of your own with \emph{le portique}
  \item \emph{Recall:} say the French for the customs officer, then the French for to declare, then say \emph{le portique} again
\end{itemize}

\begin{tcolorbox}[breakable,colback=teal!4,colframe=teal!35!black,title={Before you move on}]
What does le portique mean? (\textquotedblleft{}The scanner gate (at security); the portico\textquotedblright{}.) Say it once more.
\end{tcolorbox}
\section[fouiller]{fouiller — to search (a bag, a person)}

\label{lesson:FR-C417-fouiller}



\begin{quote}
Before the new one: say the French for to declare, then the French for the scanner gate.
\end{quote}

\subsection*{You'll want to know: fouiller}
\textbf{fouiller} — \textquotedblleft{}to search (a bag, a person)\textquotedblright{}.

\begin{quote}
Un agent a fouillé mon sac à dos.
\end{quote}

\begin{quote}
\emph{An officer searched my rucksack.}
\end{quote}

\emph{La fouille} — the search.

\begin{cousinweb}
\emph{fouiller} goes back to Vulgar Latin \emph{fodiculāre}, from Latin \emph{fodere} 'to dig' (Lewis \& Short); EtymDB links it with \emph{fouir} 'to dig'.
\end{cousinweb}

\begin{grammarlens}[title={What you've built}]
One more word for telling, explaining and describing at B1.
\end{grammarlens}

\subsection*{Your turn}
\begin{itemize}
\raggedright
  \item \emph{Say it:} \emph{fouiller}
  \item \emph{Say it:} \emph{fouiller}, once more
  \item \emph{Say it:} a sentence of your own with \emph{fouiller}
  \item \emph{Recall:} say the French for to declare, then the French for the scanner gate, then say \emph{fouiller} again
\end{itemize}

\begin{tcolorbox}[breakable,colback=teal!4,colframe=teal!35!black,title={Before you move on}]
What does fouiller mean? (\textquotedblleft{}To search (a bag, a person)\textquotedblright{}.) Say it once more.
\end{tcolorbox}
\section[transiter]{transiter — to change planes, to pass through}

\label{lesson:FR-C417-transiter}



\begin{quote}
Before the new one: say the French for the scanner gate, then the French for to search.
\end{quote}

\subsection*{You'll want to know: transiter}
\textbf{transiter} — \textquotedblleft{}to change planes, to pass through\textquotedblright{}.

\begin{quote}
Ce vol est moins cher, mais il transite par deux aéroports.
\end{quote}

\begin{quote}
\emph{This flight is cheaper, but it goes through two airports.}
\end{quote}

\emph{Le transit} — transit.

\begin{cousinweb}
\emph{transiter} is built on \emph{transit}, a learned borrowing from Latin \emph{trānsitus} 'a passing across', from \emph{trānsīre}, built on \emph{īre} 'to go' (Lewis \& Short).
\end{cousinweb}

\begin{grammarlens}[title={What you've built}]
One more word for telling, explaining and describing at B1.
\end{grammarlens}

\subsection*{Your turn}
\begin{itemize}
\raggedright
  \item \emph{Say it:} \emph{transiter}
  \item \emph{Say it:} \emph{transiter}, once more
  \item \emph{Say it:} a sentence of your own with \emph{transiter}
  \item \emph{Recall:} say the French for the scanner gate, then the French for to search, then say \emph{transiter} again
\end{itemize}

\begin{tcolorbox}[breakable,colback=teal!4,colframe=teal!35!black,title={Before you move on}]
What does transiter mean? (\textquotedblleft{}To change planes, to pass through\textquotedblright{}.) Say it once more.
\end{tcolorbox}
